\chapter{Carbohydrates}

\section*{Definition and Overview}
Carbohydrates are one of the three main nutrients in food, along with protein and fat, and they are the body's preferred and most efficient source of energy. Chemically they are sugars, starches, and fibres, and they are broken down in the gut mainly into glucose, which is carried in the blood and used by cells for fuel. Carbohydrates come in two broad forms: \textbf{simple} carbohydrates, which are sugars found in fruit, honey, and refined products, and \textbf{complex} carbohydrates, which are the starches and fibre found in whole grains, legumes, and vegetables. The type and amount of carbohydrate eaten have a major influence on blood sugar levels, body weight, dental health, and long-term risk of chronic disease.

\section*{Epidemiology}
Carbohydrate intake varies widely between populations, but in most diets carbohydrates provide around half of daily energy. In many high-income countries, people eat more refined and added sugars than recommended, which contributes to rising rates of overweight, obesity, type 2 diabetes, and tooth decay. At the same time, very low-carbohydrate and ketogenic diets have become popular for weight loss and for certain medical conditions, and global consumption of fibre remains well below recommended levels in most populations. Health authorities consistently recommend that most dietary energy from carbohydrates should come from fibre-rich, unrefined sources, and that added sugars should be limited to no more than about 10\% of daily energy.

\section*{Aetiology and Pathophysiology}
Carbohydrates are not a disease, but the way the body handles them is central to health. Their effects depend on their structure and the rate at which they are digested:
\begin{itemize}
    \item \textbf{Simple sugars:} absorbed quickly, causing a fast rise in blood glucose and a correspondingly rapid insulin response
    \item \textbf{Complex starches:} digested more slowly, giving a gentler and more sustained rise in blood glucose
    \item \textbf{Fibre:} largely undigested by human enzymes; it slows the release of sugar, feeds healthy gut bacteria, and adds bulk to stools
    \item \textbf{Glycaemic index:} a measure of how quickly a food raises blood sugar, so low-glycaemic-index foods are generally more favourable
    \item \textbf{Insulin resistance:} habitual excess of refined carbohydrate and body fat can make cells less responsive to insulin, raising blood sugar and leading to type 2 diabetes
\end{itemize}
Individual tolerance to carbohydrate varies with age, activity, muscle mass, and metabolic health, which is why there is no single intake that suits everyone.

\section*{Clinical Features}
Carbohydrates affect the body in ways that may produce symptoms when intake is out of balance:
\begin{itemize}
    \item Rapid blood sugar swings can cause fatigue, shakiness, hunger, and poor concentration
    \item High intake of refined sugars is linked to weight gain, dental decay, and non-alcoholic fatty liver disease
    \item Very low carbohydrate intake may cause headache, tiredness, bad breath, and constipation, especially in the early weeks of a restrictive diet
    \item A lack of fibre causes constipation and is associated with an unhealthy gut
    \item In people with diabetes, unplanned carbohydrate intake can cause high blood sugar, while too little can cause hypoglycaemia with sweating, tremor, and confusion
    \item Certain carbohydrates (FODMAPs) and fructose can trigger bloating, wind, and pain in people with irritable bowel syndrome
\end{itemize}

\section*{Differential Diagnosis}
Symptoms such as fatigue, headaches, or bloating after meals can have other causes:
\begin{itemize}
    \item Iron deficiency, thyroid disorders, or other metabolic conditions
    \item Poor sleep, stress, or depression
    \item Food intolerance or allergy, such as lactose intolerance
    \item Irritable bowel syndrome, coeliac disease, or other gut conditions
    \item Anaemia or chronic illness causing general tiredness
    \item Medication side effects
\end{itemize}

\section*{When to Seek Medical Care}
Most people can manage their carbohydrate intake through healthy eating habits. See a doctor or dietitian if:
\begin{itemize}
    \item You are trying to lose weight or manage diabetes and need reliable, individualised guidance
    \item You are planning a restrictive or very low-carbohydrate diet and have a medical condition
    \item You experience frequent, severe bloating or gut symptoms after eating
    \item You have symptoms of high or low blood sugar, such as unusual thirst, frequent urination, fatigue, or episodes of shakiness and sweating
    \item You have diabetes and are unsure how to match insulin or medicines to your meals
\end{itemize}
\textbf{Call 000 (or local emergency services) for:}
\begin{itemize}
    \item Severe confusion, drowsiness, or seizures from very low blood sugar
    \item Signs of diabetic ketoacidosis: severe thirst, vomiting, abdominal pain, rapid breathing, and a fruity breath smell
    \item Collapse or loss of consciousness
\end{itemize}

\section*{Diagnosis and Investigations}
Carbohydrate handling is assessed mainly through blood tests and dietary review rather than by testing the diet itself:
\begin{itemize}
    \item \textbf{Dietary history:} a structured review of eating patterns with a doctor or dietitian
    \item \textbf{Fasting blood glucose:} a screening test for diabetes and prediabetes
    \item \textbf{HbA1c:} a measure of average blood sugar control over the previous three months
    \item \textbf{Oral glucose tolerance test:} measures the blood sugar response to a standard glucose drink
    \item \textbf{Lipid profile:} cholesterol and triglycerides, which are influenced by carbohydrate and fat intake
    \item \textbf{Gut investigations:} coeliac blood tests, hydrogen breath tests, or a supervised trial of a low-FODMAP diet when gut symptoms are prominent
\end{itemize}

\section*{Management}
There is no single perfect way to eat carbohydrate, but general principles are well supported by evidence:
\begin{itemize}
    \item \textbf{Choose mostly whole, minimally processed carbohydrates:} whole grains, legumes, fruit, and vegetables
    \item \textbf{Limit added sugars and sugary drinks:} including soft drinks, fruit juice, and sweet treats
    \item \textbf{Include plenty of fibre:} aim for a variety of plant foods each day, increasing intake gradually
    \item \textbf{Match portions to activity and body size} rather than following fixed rules
    \item \textbf{For people with diabetes:} carbohydrate counting and spacing of meals, arranged with a dietitian or diabetes educator
    \item \textbf{For weight loss:} a modest reduction in total energy, often by trimming refined carbohydrates and adding protein and vegetables
    \item \textbf{For irritable bowel syndrome:} a temporary low-FODMAP diet under professional guidance
    \item \textbf{Therapeutic low-carbohydrate diets:} used under medical supervision for some people with epilepsy or specific metabolic conditions
\end{itemize}

\section*{Complications}
\begin{itemize}
    \item Obesity and its consequences from a long-term excess of refined carbohydrates and total energy
    \item Type 2 diabetes and prediabetes
    \item Dental decay from frequent exposure to sugars
    \item Non-alcoholic fatty liver disease
    \item Constipation and diverticular disease from too little fibre
    \item Nutrient shortfalls, constipation, or disordered eating from overly restrictive diets
    \item Difficulty managing blood sugar levels in people with diabetes
\end{itemize}

\section*{Prognosis and Outcomes}
For most healthy people, a balanced intake of carbohydrates is associated with good long-term health, and higher fibre intake lowers the risk of heart disease, bowel cancer, and type 2 diabetes. People who adopt very low-carbohydrate diets often lose weight in the short term, but long-term weight maintenance depends on whether the pattern is sustainable and fits the individual's preferences. With appropriate medical and dietetic support, people with diabetes and other metabolic conditions can control their blood sugar well while still eating a satisfying and nutritious diet.

\section*{Prevention and Self-Care}
\begin{itemize}
    \item Base meals on vegetables, fruits, legumes, and whole grains
    \item Read food labels for total and added sugars
    \item Drink water instead of sugary drinks
    \item Snack on fruit, nuts, or yoghurt rather than processed sweet foods
    \item Include beans, lentils, and other high-fibre foods regularly
    \item Be physically active, which helps the body use carbohydrate efficiently
    \item If you have diabetes or a gut condition, work with a dietitian rather than following self-prescribed diets
    \item Make changes gradually so that they are sustainable in the long term
\end{itemize}

\section*{Sources and Further Reading}
The following sources provide authoritative, up-to-date information on carbohydrates and healthy eating.
\begin{itemize}
    \item NHS. \textit{Eat well: starchy foods and sugar in your diet.} https://www.nhs.uk/live-well/eat-well/
    \item World Health Organization. \textit{Carbohydrate intake and healthy diet guidance.} https://www.who.int/
    \item Centers for Disease Control and Prevention. \textit{Carbohydrates and diabetes.} https://www.cdc.gov/diabetes/
    \item MedlinePlus. \textit{Carbohydrates in your diet.} https://medlineplus.gov/carbohydrates.html
    \item Mayo Clinic. \textit{Carbohydrates: how carbs fit into a healthy diet.} https://www.mayoclinic.org/healthy-lifestyle/nutrition-and-healthy-eating/
\end{itemize}
